% /Users/pikachu/books/payments-book/src/parts/part3/ch18-open-banking.tex
\chapter{Open Banking, A2A и ISO 20022}
\label{ch:open-banking}
\label{ch:global}

\begin{kaobox}[title=Что вы узнаете из этой главы]
  Зачем PSD2 и~SCA открыли банкам путь к~модели доступа через API
  к~счёту; чем AISP и~PISP отличаются от~карточных участников;
  почему платежи счёт-счёт (account-to-account, A2A)~--- это
  отдельный rail, а~не~ещё одна кнопка checkout; и~как ISO~20022
  становится общим языком для~банковских сообщений, Open Banking
  и~трансграничных переводов.
\end{kaobox}

Представим checkout, в~котором покупатель не~вводит PAN и~не~переходит
в~карточный challenge flow, а~подтверждает платёж в~приложении своего
банка. На~поверхности это выглядит как~ещё один способ оплаты.
Но~инженерно здесь меняется сама модель: центром системы становится
не~карта, а~счёт, доступ к~нему и~набор договорённостей на~уровне API.
Эта глава отвечает на~вопрос, как такая модель устроена на~уровне
регуляторики, платёжного rail и~языка сообщений.


\section{PSD2 и SCA: зачем банку открывать доступ к~счёту}
\label{sec:psd2-sca}

PSD2 изменила европейский платёжный ландшафт не~тем, что придумала
новый интерфейс, а~тем, что перестроила саму рамку доступа к~счёту.
Директива ввела лицензируемых третьих участников и~обязала банки
открывать им доступ к~данным и~инициированию платежей при~наличии
согласия клиента~\sidecite{psd2}. Иначе говоря, то, что раньше было
частным банковским каналом, стало регулируемым интерфейсом.

Одновременно с~этим европейский регулятор сформулировал требование
усиленной аутентификации (Strong Customer Authentication). Идея проста:
если доступ к~счёту
открывается через электронные каналы, банк должен убедиться, что
плательщик подтверждает операцию не~одним, а~как минимум двумя
независимыми факторами из~трёх классов: знание, владение,
биометрия~\sidecite{eba-rts-sca}. Для~чисто продуктового мышления
это выглядит как~трение; для~инженерной модели~--- как способ
сделать открытый доступ управляемым.

При~этом SCA никогда не~задумывалась как~механический шаг для~каждой
операции. Регулятор допускает исключения: low value, trusted beneficiary,
merchant-initiated и~TRA (Transaction Risk Analysis). Эти исключения
важны не~сами по~себе, а~как граница между безопасностью и~конверсией.
Без~них Open Banking быстро превратился~бы в~канал, формально открытый,
но~практически слишком тяжёлый для~повседневного использования.
\sidecite{eba-rts-sca}

\begin{important}
  Исключение из~SCA~--- не~право мерчанта и~не~гарантия бесшовного пути.
  Мерчант, эквайер или~PISP может только запросить послабление,
  но~окончательное решение принимает банк плательщика. Поэтому
  продуктовая логика \enquote{мы сейчас просто попросим exemption}
  всегда неполна: с~точки зрения системы это лишь приглашение
  к~упрощённому сценарию, а~не~обязательство его~дать.
\end{important}

Локальный вывод таков. PSD2 и~SCA не~создают платёжный rail сами
по~себе, но~задают режим, в~котором доступ к~банковскому счёту
становится машиночитаемым, лицензируемым и~проверяемым. Именно
в~этом пространстве потом появляются интерфейсы Open Banking
и~A2A-платежи.


\section{Open Banking: AISP, PISP и~поток согласия}
\label{sec:open-banking}

Как только счёт становится доступным через регулируемый интерфейс,
возникают два новых типа участников. AISP (Account Information Service
Provider) читает данные счёта и~историю транзакций, а~PISP (Payment
Initiation Service Provider) инициирует платёж со~счёта клиента
в~пользу мерчанта или~другого получателя~\sidecite{psd2}. Важно
зафиксировать различие: AISP работает с~информацией, PISP~--- с~самой
операцией. В~карточном мире эти функции часто скрыты за~интерфейсом PSP;
в~Open Banking они разведены явно и~юридически.

Возникает технический вопрос: как сделать так, чтобы один и~тот~же
продукт не~интегрировался заново с~каждым банком? Ответом стали
отраслевые API-стандарты.

В~Великобритании эту роль играет стандарт Open Banking UK,
а~в~континентальной Европе~--- профиль Berlin Group
NextGenPSD2~\sidecite{openbanking-uk}\sidecite{berlin-group}.

Оба подхода описывают один и~тот~же поток из~трёх шагов:
согласие клиента, выдача токена доступа, обращение к~API счёта
или~инициирования платежа.

\todofig{PISP flow. Участники: клиент, мерчант/PISP, банк клиента.
  Шаги: выбор оплаты со счёта; редирект в банк; consent + SCA;
  выдача токена и возврат клиента; платёж через API; статус
  исполнения; подтверждение заказа.}
% \label{fig:pisp-flow}

В~этой схеме поток согласия играет ту~же роль, что и~сбор карточных
данных в~классическом checkout: он~определяет, на~каких правах
третья сторона вообще входит в~операцию. Но~в~отличие от~карточного
мира здесь авторизация пользователя и~управление согласием становятся
частью банковского интерфейса, а~не~надстройкой PSP. Поэтому
Open Banking лучше понимать не~как \enquote{ещё один API банка},
а~как новую точку разграничения ответственности между банком,
агрегатором и~мерчантом.

В~России разговор об~Open API идёт уже несколько лет, но~картина
остаётся менее унифицированной, чем~в~европейском
контуре~\todocheck{актуальный статус обязательного стандарта
  Open API Банка России на~начало 2026~года}. Сам по~себе этот
разрыв показателен: нормативная рамка определяет не~только то,
какие API существуют, но~и~то, насколько отрасль вообще может
строить поверх них массовые платёжные сценарии.


\section{A2A-платежи: отдельный rail, а~не~кнопка на~странице}
\label{sec:a2a-payments}

Когда PISP инициирует перевод со~счёта клиента, возникает следующий
вопрос: по~какому именно rail пойдут деньги? Ответом становятся
платежи счёт-счёт~--- account-to-account, или A2A. Их~важно понимать
не~как новый пользовательский интерфейс, а~как отдельную расчётную
модель, в~которой деньги движутся между банковскими счетами без
карточного посредника.

Разные юрисдикции решают эту задачу по-разному. Бразильский PIX
сделал мгновенный банковский перевод массовым розничным инструментом
повседневной оплаты~\sidecite{pix}. Индийский UPI показал другую
конструкцию: единый инфраструктурный слой и~конкуренцию приложений
поверх него~\sidecite{upi}. В~европейском контуре роль эталона
скорее выполняют мгновенные схемы SEPA, где акцент делается на~скорости
межбанковского перевода и~на~стандартизации сообщений, а~не~на~карточной
экосистеме~\sidecite{sepa-instant}. На~этом фоне СБП, разобранная
в~предыдущей главе, выглядит не~локальной экзотикой, а~частью
общего мирового сдвига в~сторону банковских instant rails.

Именно здесь становится ясно, почему A2A нельзя описывать словами
\enquote{ещё один способ checkout}. В~карточном мире у~нас есть
авторизация, отдельный клиринг, chargeback и~многоуровневая модель
ответственности. В~мире A2A-платежей логика иная. Подтверждение пользователя
и~движение денег оказываются ближе друг к~другу по~времени, операция
быстрее становится окончательной, а~принудительный механизм спора
уровня chargeback либо отсутствует, либо устроен принципиально иначе.
То~есть меняется не~кнопка на~странице, а~сами гарантии, на~которые
может рассчитывать мерчант и~покупатель.

Из~этого вырастает и~основное ограничение A2A. Такие rails сильны там,
где важны низкая стоимость, мгновенность и~прямой доступ к~банковскому
счёту. Но~карты по-прежнему выигрывают там, где нужен глобальный
охват, единый слой защиты покупателя и~отлаженная система
оспаривания. Поэтому рост A2A~--- не~конец карточного мира,
а~перераспределение задач между rails.


\section{ISO~20022 вне карточного мира}
\label{sec:iso20022}

Как только платёжный разговор выходит за~границы карточных сообщений,
почти сразу появляется ISO~20022. Этот стандарт важен не~тем, что
заменяет одно имя поля другим, а~тем, что вводит более богатую модель
данных для~финансового сообщения. В~ней у~элемента есть не~только
позиция в~строке, но~и~заранее определённая семантика, тип
и~место в~общей структуре~\sidecite{iso-20022}. Именно поэтому
ISO~20022 оказывается удобным языком там, где нужно одновременно
решать задачу маршрутизации, AML-скрининга, reconciliation
и~отчётности.

Для~разработчика особенно важны три семейства сообщений. \texttt{pain}
описывает инициирование платежа клиентом; \texttt{pacs}~--- клиринг
и~расчёты; \texttt{camt}~--- статусы, выписки и~другие сообщения
управления ликвидностью~\sidecite{iso-20022}. Если смотреть
на~Open Banking и~A2A вместе, становится видно, как эти семейства
складываются в~единый слой: API даёт удобную оболочку для~продукта,
а~ISO~20022 даёт машиночитаемую форму для~самого банковского
сообщения.

Эта роль особенно заметна в~трансграничных переводах. Для~SWIFT
переход к~CBPR+ означает, что MX становится основным режимом обмена
сообщениями после 22~ноября 2025~года, но~операционная картина не~сводится
к~простому \enquote{MT выключили, MX включили}. Остались исключения,
остаточные MT-потоки и~аварийные сценарии, которые нужно читать
не~по~общему лозунгу о~миграции, а~по~конкретным правилам и~бюллетеням
SWIFT~\todocheck{точный operational status SWIFT CBPR+
  после 22~ноября 2025~года: остаточные MT-потоки и~аварийные
  сценарии}. Практический смысл для~инженера в~том, что после
основной миграции уже недостаточно знать только \enquote{формат
  сообщения}; нужно понимать, какой профиль, какой переходный сценарий
и~какая исключительная ветка действуют для~конкретного потока.

СПФС движется в~том~же направлении: поддержка форматов MX делает
ISO~20022 важным не~только для~контура SWIFT, но~и~для~российской
межбанковской инфраструктуры~\sidecite{cbr-spfs}. Отсюда и~общий
принцип: чем дальше система уходит от~карты и~чем ближе подходит
к~банковскому счёту, тем важнее становится не~пластик и~не~PAN,
а~богатая структура самого сообщения.

\begin{practice}
  При~проектировании систем под ISO~20022 полезно заранее развести
  три слоя. Первый~--- внешний API, через который продукт разговаривает
  с~клиентом или~мерчантом. Второй~--- каноническая внутренняя модель
  платежа. Третий~--- профиль ISO~20022 или~другого межбанковского
  стандарта, в~который эта модель сериализуется. Когда эти слои
  склеены, любая миграция на~новый профиль CBPR+, новую мгновенную схему
  или~новый банковский API превращается в~переписывание всего контура
  сразу.
\end{practice}


\section*{Итоги главы}

\begin{itemize}
  \item PSD2 и~SCA открывают доступ к~счёту не~хаотически, а~через
        регулируемый набор ролей, согласий и~правил аутентификации.
  \item Open Banking разводит функции чтения данных и~инициирования
        платежа между AISP и~PISP и~строит вокруг них поток согласия
        и~стандарты API.
  \item A2A-платежи~--- это отдельный rail со~своей логикой
        окончательности, ответственности и~защиты покупателя,
        а~не~просто новый способ нарисовать checkout.
  \item ISO~20022 становится общим языком для~Open Banking,
        instant rails и~межбанковских переводов, потому что несёт
        более богатую и~машиночитаемую структуру сообщения.
\end{itemize}
